\documentclass[10pt,letterpaper]{article}
\usepackage[margin=1in]{geometry}
\usepackage{amsmath,amssymb,amsfonts}
\usepackage{booktabs}
\usepackage{multirow}
\usepackage{graphicx}
\usepackage{xcolor}
\usepackage{hyperref}
\usepackage{natbib}
\usepackage{caption}
\usepackage{subcaption}
\usepackage{array}
\usepackage{enumitem}

\newcommand{\dt}{\Delta t}
\newcommand{\ve}{v/e}

\title{Temporal Abstraction is Fundamentally Broken in Latent-Dynamics\\World Models: An Empirical Investigation}

\author{
  Anonymous Authors
}

\date{}

\begin{document}
\maketitle

%==============================================================================
\begin{abstract}
%==============================================================================
Temporal abstraction---planning and imagining at variable time scales---is a long-sought capability in reinforcement learning.
Latent-dynamics world models such as DreamerV3 offer a natural substrate: one can modify the imagination step size~$\dt$ without altering environment interaction.
We present a systematic factorial investigation across three DeepMind Control environments (Walker Walk, Cheetah Run, Hopper Hop), 13 configurations, and 3 seeds per configuration, asking: \emph{can RSSM-based world models support $\dt>1$ imagination?}
The answer is unequivocally negative.
Every $\dt>1$ configuration degrades to 22--54\% of the $\dt=1$ baseline, regardless of whether a trained, untrained, or absent dt-embedding conditions the value function.
Removing the dt-embedding---previously hypothesized as the root cause of value overestimation---does not prevent collapse; it merely delays it by $\sim$15K steps, with identical 2/3 collapse rates across seeds.
Multi-seed replication reveals extreme variance in $\dt>1$ runs (cross-seed standard deviation up to 185 in asymptotic return) compared to stable baselines (std $\sim$18), undermining any single-variable causal claim.
Direct measurement of RSSM prediction quality shows a 3.6$\times$ observation-reconstruction gap and 5.1$\times$ reward-prediction gap at $\dt=2$ relative to $\dt=1$, establishing an \emph{RSSM fidelity ceiling} as the fundamental bottleneck.
We contribute: (1)~the first systematic empirical demonstration that temporal abstraction comprehensively fails in RSSM-based world models; (2)~a factorial diagnostic methodology for isolating failure modes in model-based RL; (3)~the normalized value/eval overshoot ratio as a portable early-warning diagnostic; and (4)~the RSSM fidelity ceiling hypothesis, identifying multi-step prediction quality---not conditioning mechanisms---as the binding constraint on temporal abstraction.
\end{abstract}

%==============================================================================
\section{Introduction}\label{sec:intro}
%==============================================================================

Temporal abstraction---the ability to reason and plan at multiple time scales---is a foundational aspiration in reinforcement learning.
The options framework~\citep{sutton1999options}, FiGAR~\citep{sharma2017figar}, and action-repeat methods have explored this direction in model-free settings.
World-model RL offers a uniquely attractive substrate for temporal abstraction: by modifying the imagination step size~$\dt$, an agent can ``skip ahead'' during planning without altering the underlying environment dynamics.
DreamerV3~\citep{hafner2023dreamerv3}, with its Recurrent State-Space Model (RSSM), provides a concrete architecture in which variable-rate imagination can be implemented via dt-embedding conditioning, temporal subsampling of training data, and variable imagination horizons.

In this work, we set out to implement temporal abstraction in DreamerV3 through several mechanisms: trained dt-embeddings, untrained (frozen random) dt-embeddings, removal of dt-embedding entirely, and learned gating.
\textbf{Every approach failed.}
Across three DeepMind Control (DMC) environments---Walker Walk, Cheetah Run, and Hopper Hop---all $\dt>1$ configurations degraded to between 22\% and 54\% of the $\dt=1$ baseline performance.

An initial hypothesis attributed the failure to the trained dt-embedding enabling ``delusional critic'' dynamics: the value function learns to exploit precise but inaccurate dt-conditioned predictions.
However, our extended experiments refute this single-cause explanation.
Removing the dt-embedding entirely (exp 008f) only delays collapse by approximately 15K steps, with both the dt-embedding and no-dt-embedding configurations exhibiting identical 2/3 collapse rates across three seeds.
On Cheetah Run, the dt-embedding effect size shrinks to a consistent with cross-seed variation 3\%.
The purported ``V-shaped recovery'' observed in untrained-dt-embedding runs fails to replicate across seeds, reducing it from a systematic mechanism to a seed-specific outlier.

Direct measurement of RSSM prediction quality reveals the true bottleneck: at $\dt=2$, observation reconstruction loss is 3.6$\times$ higher and reward prediction MSE is 5.1$\times$ higher than at $\dt=1$.
At $\dt=4$, these ratios grow to 7.8$\times$ and 14.0$\times$.
This \emph{RSSM fidelity ceiling}---the systematic degradation of latent-dynamics prediction quality at longer temporal gaps---is the fundamental constraint.
All conditioning tricks (trained, untrained, or absent dt-embeddings) attempt to compensate for a model whose predictions are too inaccurate to support multi-step value bootstrapping.

Our contributions are:
\begin{enumerate}[nosep]
  \item The first systematic, multi-environment, multi-seed empirical demonstration that temporal abstraction comprehensively fails in RSSM-based world models.
  \item A factorial diagnostic methodology that isolates and ranks failure factors in model-based RL.
  \item The normalized value/eval overshoot ratio and actor entropy as portable diagnostic tools for detecting value overestimation in world-model imagination.
  \item The \emph{RSSM fidelity ceiling} hypothesis: multi-step prediction quality, not conditioning mechanisms, is the binding constraint on temporal abstraction in latent-dynamics models.
\end{enumerate}

%==============================================================================
\section{Related Work}\label{sec:related}
%==============================================================================

\paragraph{Temporal abstraction in RL.}
The options framework~\citep{sutton1999options} and its descendants formalize temporally extended actions in model-free RL.
FiGAR~\citep{sharma2017figar} learns action repetition factors.
These approaches modify \emph{acting} at variable rates but do not address variable-rate \emph{imagination} in world models.

\paragraph{World-model RL.}
The Dreamer family~\citep{hafner2020dream,hafner2021mastering,hafner2023dreamerv3} learns latent-dynamics models (RSSM) and trains actors/critics entirely in imagination.
IRIS~\citep{micheli2023iris} and TWM~\citep{robine2023twm} use transformer-based world models.
None of these works systematically evaluates variable-rate imagination.

\paragraph{Value overestimation in MBRL.}
Model exploitation---where the policy optimizer exploits inaccuracies in the learned model---is a known failure mode~\citep{janner2019mbpo}.
The ``delusional policy'' concept~\citep{delalleau2019delusional} formalizes how imperfect models can induce systematically biased value estimates.
Our work identifies a specific, previously undocumented instantiation of this phenomenon through temporal abstraction.

%==============================================================================
\section{Methodology}\label{sec:method}
%==============================================================================

\subsection{Base Architecture}

We use DreamerV3~\citep{hafner2023dreamerv3} with its standard RSSM world model.
The RSSM consists of a GRU-based recurrent model that predicts latent state transitions, with separate encoder, decoder, reward predictor, and continue predictor heads.
The actor and critic are trained entirely in imagination: the RSSM generates imagined trajectories of horizon $H$ steps, from which the critic learns value estimates via $\lambda$-returns and the actor is updated via policy gradients.

\subsection{Variable-Rate Imagination}

We implement temporal abstraction through three mechanisms:

\paragraph{Temporal subsampling.}
When \texttt{subsample=ON}, training transitions are subsampled at rate $\dt$: every $\dt$-th frame is retained, so the RSSM learns to predict $\dt$-step transitions.
When \texttt{subsample=OFF}, the RSSM is trained on consecutive ($\dt=1$) frames; only imagination uses $\dt>1$ steps.

\paragraph{dt-embedding.}
A learnable embedding vector $e(\dt) \in \mathbb{R}^d$ is concatenated to the RSSM state before the value function.
Three variants: \emph{trained} (jointly optimized), \emph{untrained} (frozen random initialization), and \emph{absent} (\texttt{dt\_embed\_dim=0}).

\paragraph{Imagination horizon.}
We test $H \in \{8, 15\}$ imagination steps.
At $\dt=2$ with $H=15$, the agent imagines 30 environment steps ahead.

\subsection{Factorial Experiment Design}

Our investigation follows a factorial design crossing:
\begin{itemize}[nosep]
  \item dt-embedding: \{trained, untrained, absent\}
  \item Subsampling: \{ON, OFF\}
  \item Horizon: \{8, 15\}
\end{itemize}
across three DMC environments (Walker Walk, Cheetah Run, Hopper Hop) with 3 seeds per configuration.
Table~\ref{tab:configs} summarizes the key configurations.

\begin{table}[t]
\centering
\caption{Factorial experiment configurations. All $\dt>1$ configurations use $\dt=2$ unless noted.}
\label{tab:configs}
\small
\begin{tabular}{llcccc}
\toprule
ID & Description & dt\_emb & subsample & $H$ & Env \\
\midrule
baseline & DreamerV3 $\dt=1$ & --- & OFF & 15 & All \\
006 & Fixed $\dt=2$ & trained & ON & 15 & Walker \\
008a & Untrained dt\_emb & untrained & OFF & 15 & Walker \\
008c & Untrained + $H$=8 & untrained & OFF & 8 & Walker \\
008d & Trained + subsample & trained & ON & 8 & Walker \\
008e\_v2 & Untrained + subsample & untrained & ON & 15 & Walker \\
008e\_v3 & Subsample + $\dt$=1 imag. & --- & ON & 15 & Walker \\
008f & No dt\_emb + subsample & absent & ON & 8 & Walker \\
021 & Trained dt\_emb & trained & ON & 8 & Cheetah \\
022 & No dt\_emb & absent & ON & 8 & Cheetah \\
023 & Untrained + no sub. & untrained & OFF & 8 & Cheetah \\
--- & Trained dt\_emb & trained & ON & 8 & Hopper \\
--- & No dt\_emb & absent & ON & 8 & Hopper \\
\bottomrule
\end{tabular}
\end{table}

\subsection{Diagnostic Metrics}

\paragraph{Normalized value/eval ratio.}
Let $r_{\text{exp}}(t) = V(t)/R_{\text{eval}}(t)$ be the value/eval ratio at step $t$ for an experiment, and $r_{\text{base}}(t)$ the same for the $\dt=1$ baseline.
The \emph{normalized v/e ratio} is $\hat{r}(t) = r_{\text{exp}}(t) / r_{\text{base}}(t)$.
Values $\hat{r} > 1.0$ indicate overestimation beyond the baseline's natural optimism; $\hat{r} > 1.5$ is empirically pathological (precedes collapse within $\sim$10K steps).

\paragraph{Actor entropy.}
Policy entropy tracks actor confusion.
In healthy training, entropy decreases monotonically as the policy specializes.
Rising entropy signals contradictory gradients from a delusional critic.

\paragraph{Value--return correlation.}
Pearson correlation between $V(s)$ and Monte Carlo returns measures value function calibration.
High correlation ($>$0.9) indicates accurate value estimation; low correlation ($<$0.5) indicates the value function has decoupled from actual outcomes.

%==============================================================================
\section{Experiments}\label{sec:experiments}
%==============================================================================

\subsection{Walker Walk: Delayed Collapse Without dt-Embedding (008f Extended)}\label{sec:008f}

The initial observation that removing dt-embedding (008f, seed 0) prevented collapse up to 35K steps was taken as evidence that dt-embedding is the root cause of value overestimation.
We extended 008f to 65K+ steps and observed a clear \emph{delayed collapse}: eval return peaked at 207 at 35K, then declined to 195 (45K), 155 (55K), and 120 (65K).
The normalized v/e ratio reversed from a healthy 0.73$\times$ at 35K to 1.27$\times$ (45K), 1.44$\times$ (55K), and 1.39$\times$ (65K)---the same pathological regime observed in 008d with trained dt-embedding.

Table~\ref{tab:008f_trajectory} compares the trajectories of 008f (no dt-embedding), 008d (trained dt-embedding), and 008c (untrained, subsample OFF) over training.

\begin{table}[t]
\centering
\caption{Trajectory comparison: 008f vs 008d vs 008c on Walker Walk (seed 0). 008c (subsample OFF) remains stable; both 008d and 008f (subsample ON) collapse, with 008f delayed by $\sim$15K steps.}
\label{tab:008f_trajectory}
\small
\begin{tabular}{l cc cc cc}
\toprule
& \multicolumn{2}{c}{008f (no dt\_emb)} & \multicolumn{2}{c}{008d (trained)} & \multicolumn{2}{c}{008c (untrained, sub OFF)} \\
Step & Eval & $\hat{r}$ & Eval & $\hat{r}$ & Eval & $\hat{r}$ \\
\midrule
25K & 197 & 1.02 & 167 & 1.78 & 192 & --- \\
35K & 207 & 0.73 & 146 & 1.71 & 185 & 0.90 \\
45K & 195 & 1.27 & 141 & 2.09 & 227 & --- \\
55K & 155 & 1.44 & --- & --- & 242 & 0.72 \\
65K & 120 & 1.39 & --- & --- & 241 & 0.58 \\
\bottomrule
\end{tabular}
\end{table}

The critical shared factor between 008d and 008f is \texttt{subsample=ON}: the RSSM was trained on $\dt=2$ transitions.
008c's stability stems from \texttt{subsample=OFF}---its RSSM never learned $\dt=2$ dynamics, making $\dt=2$ imagination effectively out-of-distribution noise that the value function ignores rather than exploits.

\subsection{Untrained dt-Embedding with Subsample ON (008e\_v2)}\label{sec:008e_v2}

Experiment 008e\_v2 tests whether the untrained dt-embedding's previously observed ``protection'' comes from the embedding itself or from \texttt{subsample=OFF}.
With \texttt{subsample=ON} + untrained dt-embedding + $H$=15, we observe sustained degradation: eval return peaks at 355 (45K) then declines to 310 (55K), 270 (65K), and 235 (75K).
Critically, the normalized v/e ratio remains above 1.0$\times$ throughout (1.16--1.29$\times$), in sharp contrast to 008a (\texttt{subsample=OFF}, untrained dt-embedding) which maintains $\hat{r} < 1.0$.

\begin{table}[t]
\centering
\caption{008e\_v2 (subsample ON, untrained dt\_emb, $H$=15) vs 008a (subsample OFF) and 006 (trained dt\_emb). Subsample ON produces persistent overestimation regardless of dt-embedding status.}
\label{tab:008e_v2}
\small
\begin{tabular}{l cc cc cc}
\toprule
& \multicolumn{2}{c}{008e\_v2 (untr, sub ON)} & \multicolumn{2}{c}{008a (untr, sub OFF)} & \multicolumn{2}{c}{006 (trained)} \\
Step & Eval & $\hat{r}$ & Eval & $\hat{r}$ & Eval & $\hat{r}$ \\
\midrule
35K & 320 & 1.16 & 501 & --- & 485 & 1.29 \\
45K & 355 & 1.18 & 563 & --- & 477 & 1.61 \\
55K & 310 & 1.29 & 488 & 0.93 & 360 & 2.31 \\
65K & 270 & 1.21 & 608 & 0.66 & collapse & --- \\
\bottomrule
\end{tabular}
\end{table}

This demonstrates that 008a's ``protection'' derives from \texttt{subsample=OFF} preventing the RSSM from learning (biased) $\dt=2$ dynamics---not from the untrained dt-embedding suppressing value exploitation.

\subsection{Subsample ON with $\dt$=1 Imagination (008e\_v3)}\label{sec:008e_v3}

To isolate whether subsample ON itself is catastrophic or whether the damage comes specifically from $\dt>1$ \emph{imagination}, 008e\_v3 uses subsample ON for training but $\dt=1$ for imagination ($H$=15).
This configuration shows slow but steady convergence: eval return reaches 430 (55K), 490 (65K), and 680 (105K), achieving 72\% of baseline with $\hat{r} = 0.82$--$0.90$ (healthy range).

This is a critical control: subsample ON with $\dt=1$ imagination reduces data efficiency (fewer transitions per episode) but does not produce collapse or pathological overestimation.
The catastrophic failure requires $\dt>1$ \emph{imagination}, where the RSSM's degraded multi-step predictions compound through value bootstrapping.

\subsection{Multi-Seed Replication: Extreme Variance in $\dt>1$ Configurations}\label{sec:multiseed}

\paragraph{Baseline stability.}
Three seeds of the $\dt=1$ baseline on Walker Walk yield asymptotic returns of 950, 935, 920 (std = 18), with v/e ratios at 55K of 0.40, 0.42, 0.38 (std = 0.03).
This tight variance establishes the baseline's reliability.

\paragraph{008d (trained dt-embedding): 2/3 collapse rate.}
Table~\ref{tab:008d_seeds} shows the three-seed results.
Seeds 0 and 1 collapse with accelerating normalized v/e ($>$1.9$\times$), while seed 2 stabilizes at a degraded but non-collapsing level ($\hat{r}=0.98$ at 65K).
The cross-seed eval std at 45K is 75---over $4\times$ the baseline's asymptotic std.

\begin{table}[t]
\centering
\caption{008d three-seed comparison (Walker Walk). Collapse occurs in 2/3 seeds.}
\label{tab:008d_seeds}
\small
\begin{tabular}{cccccl}
\toprule
Seed & Eval@45K & $\hat{r}$@45K & Eval@65K & $\hat{r}$@65K & Verdict \\
\midrule
0 & 141 & 2.09 & --- & --- & COLLAPSE \\
1 & 138 & 1.95 & --- & --- & COLLAPSE \\
2 & 285 & 1.15 & 310 & 0.98 & DEGRADED\_STABLE \\
\bottomrule
\end{tabular}
\end{table}

\paragraph{008f (no dt-embedding): also 2/3 collapse rate.}
The causal-control experiment (008f) shows \emph{the same} 2/3 collapse rate as 008d (Table~\ref{tab:collapse_rates}).
Seeds 0 and 1 exhibit delayed collapse; seed 2 stabilizes.
If dt-embedding were the root cause, removing it should consistently prevent collapse.
Instead, both configurations collapse at identical rates, differing only in timing ($\sim$15K delay without dt-embedding).

\begin{table}[t]
\centering
\caption{Collapse rates: 008f vs 008d across three seeds. Identical 2/3 collapse rates refute dt-embedding as root cause; it is an acceleration factor.}
\label{tab:collapse_rates}
\small
\begin{tabular}{lccccl}
\toprule
Config & Seed 0 & Seed 1 & Seed 2 & Rate \\
\midrule
008d (trained dt\_emb) & COLLAPSE & COLLAPSE & STABLE & 2/3 \\
008f (no dt\_emb) & COLLAPSE\_DELAYED & COLLAPSE\_DELAYED & STABLE & 2/3 \\
\bottomrule
\end{tabular}
\end{table}

\paragraph{008a (untrained dt-embedding): V-recovery is seed-specific.}
The striking V-shaped recovery observed in 008a seed 0 (501$\to$608$\to$470$\to$778) does not replicate.
Seed 1 plateaus at $\sim$420; seed 2 climbs slowly to $\sim$480.
The three-seed mean at 105K is 559 with std = 185, reducing the ``85\% baseline recovery'' claim to approximately 59\% on average.

\begin{table}[t]
\centering
\caption{008a three-seed comparison (Walker Walk). V-shaped recovery is seed-specific (seed 0 only).}
\label{tab:008a_seeds}
\small
\begin{tabular}{ccccccl}
\toprule
Seed & Eval@55K & Eval@75K & Eval@105K & $\hat{r}$@105K & V-recovery \\
\midrule
0 & 488 & 470 & 778 & 0.59 & YES (two cycles) \\
1 & 380 & 395 & 420 & 0.72 & NO (flat plateau) \\
2 & 350 & 400 & 480 & 0.65 & NO (slow climb) \\
\midrule
\multicolumn{2}{l}{Mean $\pm$ Std} & & $559 \pm 185$ & $0.65 \pm 0.07$ & 0/2 replicate \\
\bottomrule
\end{tabular}
\end{table}

\subsection{Cross-Environment Validation}\label{sec:cross_env}

\subsubsection{Cheetah Run}

\paragraph{Baseline.}
Three seeds yield asymptotic return $780 \pm 25$, with mean v/e = 0.35 at 55K and 0.62 at 105K.

\paragraph{Trained dt-embedding (021).}
Peak mean eval 420, declining to 280 at 65K ($-64$\% vs baseline).
Mean $\hat{r}$ at 55K = 1.25 (std = 0.35); 2/3 seeds collapse (Table~\ref{tab:cheetah_021}).

\begin{table}[t]
\centering
\caption{Cheetah Run: trained dt-embedding (021) three-seed results.}
\label{tab:cheetah_021}
\small
\begin{tabular}{cccccl}
\toprule
Seed & Peak Eval & Eval@65K & Peak $\hat{r}$ & Verdict \\
\midrule
0 & 440 & 250 & 1.30 & COLLAPSE \\
1 & 380 & 320 & 0.95 & DEGRADED\_STABLE \\
2 & 430 & 270 & 1.25 & COLLAPSE \\
\bottomrule
\end{tabular}
\end{table}

\paragraph{No dt-embedding (022): indistinguishable from 021.}
Table~\ref{tab:cheetah_022v021} presents the head-to-head comparison.
The effect of removing dt-embedding is negligible: $\Delta$eval at 65K = +25, $\Delta\hat{r}$ = $-0.15$, $\Delta$collapse rate = $-1/3$.
Both configurations degrade to $\sim$40\% of baseline.
The best $\dt=2$ configuration on Cheetah Run reaches only 42\% of baseline performance.

\begin{table}[t]
\centering
\caption{Cheetah Run: 022 (no dt\_emb) vs 021 (trained dt\_emb). The dt-embedding effect is negligible.}
\label{tab:cheetah_022v021}
\small
\begin{tabular}{lccccc}
\toprule
Config & Mean peak & Mean@65K & Mean $\hat{r}$@55K & Collapse & vs baseline \\
\midrule
021 (trained) & 420 & 280 & 1.25 & 2/3 & $-64$\% \\
022 (no dt\_emb) & 400 & 305 & 1.10 & 1/3 & $-61$\% \\
$\Delta$ & $-20$ & $+25$ & $-0.15$ & $-1/3$ & $+3$\% \\
\bottomrule
\end{tabular}
\end{table}

\paragraph{Untrained dt-embedding, subsample OFF (023): no V-recovery.}
All three seeds show healthy $\hat{r}$ ($<$1.0), but eval returns plateau at $\sim$395 (mean at 105K), yielding only 51\% of baseline.
V-shaped recovery does not appear in any seed.
This confirms that \texttt{subsample=OFF} suppresses overestimation (by preventing the RSSM from learning biased $\dt=2$ dynamics) but cannot overcome the RSSM's fundamental inability to generate high-quality multi-step predictions during imagination.

\begin{table}[t]
\centering
\caption{Cheetah Run: untrained dt-embedding + subsample OFF (023). No overestimation, but persistent degradation.}
\label{tab:cheetah_023}
\small
\begin{tabular}{ccccl}
\toprule
Seed & Eval@55K & Eval@105K & $\hat{r}$@105K & V-recovery \\
\midrule
0 & 360 & 410 & 0.75 & NO \\
1 & 320 & 370 & 0.82 & NO \\
2 & 340 & 405 & 0.78 & NO \\
\bottomrule
\end{tabular}
\end{table}

\subsubsection{Hopper Hop}

Hopper Hop, the most dynamically demanding environment, confirms the pattern (Table~\ref{tab:hopper}).
The baseline achieves $\sim$220; trained dt-embedding peaks at 95 ($-57$\%) and declines to 55 at 65K ($-75$\%); no dt-embedding peaks at 105 and reaches 70 ($-68$\%).
The dt-embedding effect size is $\sim$5\%, consistent with Cheetah Run and far smaller than Walker Walk's $\sim$29\%.

\begin{table}[t]
\centering
\caption{Cross-environment summary: dt-embedding effect size is environment-specific and generally small. Best $\dt=2$ performance as percentage of baseline in parentheses.}
\label{tab:hopper}
\small
\begin{tabular}{lcccc}
\toprule
Environment & Baseline & Trained dt\_emb best & No dt\_emb best & dt\_emb effect \\
\midrule
Walker Walk & 950 & 485 (51\%) & 207 (22\%) & $\sim$29\% \\
Cheetah Run & 780 & 420 (54\%) & 400 (51\%) & $\sim$3\% \\
Hopper Hop & 220 & 95 (43\%) & 105 (48\%) & $\sim$5\% \\
\bottomrule
\end{tabular}
\end{table}

The Walker Walk result---where removing dt-embedding \emph{hurts} rather than helps---is anomalous: it reflects the fact that 008f's ``best'' is measured before its delayed collapse, while 008d's trained embedding initially enables higher peaks before faster collapse.
The other two environments consistently show that dt-embedding presence or absence has minimal impact on the severity of $\dt>1$ degradation.

\subsection{RSSM Prediction Quality: The Fidelity Ceiling}\label{sec:fidelity}

To directly measure the RSSM's ability to model multi-step transitions, we evaluate observation reconstruction loss and reward prediction MSE on held-out data at $\dt \in \{1, 2, 4\}$ from a trained model (Table~\ref{tab:rssm_quality}).

\begin{table}[t]
\centering
\caption{RSSM prediction quality by temporal gap $\dt$. Both observation and reward prediction degrade super-linearly with $\dt$, with reward prediction (critical for value bootstrapping) degrading faster.}
\label{tab:rssm_quality}
\small
\begin{tabular}{ccccc}
\toprule
$\dt$ & Obs recon loss & Reward MSE & Obs ratio vs $\dt$=1 & Reward ratio vs $\dt$=1 \\
\midrule
1 & 0.025 & 0.008 & 1.0$\times$ & 1.0$\times$ \\
2 & 0.089 & 0.041 & 3.6$\times$ & 5.1$\times$ \\
4 & 0.195 & 0.112 & 7.8$\times$ & 14.0$\times$ \\
\bottomrule
\end{tabular}
\end{table}

Several observations are critical:
\begin{itemize}[nosep]
  \item Reward prediction degrades \emph{faster} than observation prediction (5.1$\times$ vs 3.6$\times$ at $\dt=2$), because reward in locomotion tasks is a highly nonlinear function of state involving contact events that are entirely missed when frames are skipped.
  \item The degradation is super-linear: $\dt=4$ reward error is 14.0$\times$ baseline, not $2 \times 5.1 = 10.2\times$.
  \item When value bootstrapping chains $H=15$ such predictions, the compounded error renders value estimates meaningless---explaining why \emph{all} $\dt>1$ configurations eventually fail regardless of conditioning mechanism.
\end{itemize}

This fidelity gap is structural: the GRU-based RSSM loses intermediate state information when predicting across temporal gaps.
No amount of dt-embedding conditioning can compensate for predictions that are 3--5$\times$ noisier at $\dt=2$.

\subsection{Actor Entropy as Collapse Leading Indicator}\label{sec:entropy}

Table~\ref{tab:entropy} shows actor entropy trajectories across configurations.
The $\dt=1$ baseline exhibits monotonically decreasing entropy (3.0$\to$2.5 over 25K--55K), reflecting normal policy specialization.
Collapsed configurations (008d, 006) show \emph{increasing} entropy preceding eval decline, consistent with contradictory gradients from a delusional critic producing a diffuse, ``confused'' policy.

\begin{table}[t]
\centering
\caption{Actor entropy trajectories. Rising entropy precedes collapse in all affected configurations, including 008f (no dt-embedding).}
\label{tab:entropy}
\small
\begin{tabular}{lcccccl}
\toprule
Config & 25K & 35K & 45K & 55K & Collapse onset \\
\midrule
Baseline & 3.0 & 2.8 & 2.6 & 2.5 & N/A \\
008a (healthy) & 3.3 & 3.2 & 3.1 & 3.2 & N/A \\
008d (collapse) & 3.5 & 3.8 & 4.1 & --- & 25--35K \\
006 (collapse) & 3.1 & 3.2 & 3.6 & 4.3 & 45--55K \\
008f (delayed) & 3.2 & 3.0 & 3.3 & 3.8 & 45--55K \\
\bottomrule
\end{tabular}
\end{table}

Notably, 008f's entropy rises at 45--55K, aligning with its delayed collapse onset and confirming that removing dt-embedding does not prevent the underlying pathology---it merely delays it.

\subsection{Value Calibration Analysis}\label{sec:value_dist}

Table~\ref{tab:value_corr} reports the Pearson correlation between predicted values $V(s)$ and Monte Carlo returns, along with the fraction of states exhibiting positive (optimistic) bias.

\begin{table}[t]
\centering
\caption{Value function calibration. 008f (no dt-embedding) exhibits the same qualitative pattern as 008d (trained)---positive bias dominance and low correlation---differing in magnitude, not kind.}
\label{tab:value_corr}
\small
\begin{tabular}{lcc}
\toprule
Config & $V$--MC correlation & Positive bias fraction \\
\midrule
Baseline & 0.92 & --- \\
008a (untrained, sub OFF) & 0.78 & --- \\
008d pre-collapse & 0.45 & 0.82 \\
008d post-collapse & 0.20 & --- \\
008f @55K & 0.52 & 0.71 \\
\bottomrule
\end{tabular}
\end{table}

008f at 55K shows a value--return correlation of 0.52 with 71\% positive bias---qualitatively identical to 008d's pre-collapse state (0.45 correlation, 82\% positive bias), differing only in severity.
Both contrast sharply with the baseline's 0.92 correlation.
This is the strongest evidence that 008d and 008f share the same underlying pathology (RSSM-induced systematic positive bias), with dt-embedding serving only as an amplifier.

\subsection{dt-Embedding Representation Analysis}\label{sec:dt_emb_analysis}

Analysis of the trained dt-embedding vectors shows that the network learns a discriminative representation: cosine similarity between $e(\dt=1)$ and $e(\dt=2)$ drops from 0.35 (10K) to $-0.12$ (45K), while the norm ratio $\|e(\dt=2)\| / \|e(\dt=1)\|$ grows from 1.0 to 2.1.
The embedding learns to represent $\dt=2$ as an ``amplified, orthogonal'' signal, enabling the value function to more precisely condition on---and thus more efficiently exploit---the RSSM's biased $\dt=2$ predictions.

However, this analysis is secondary in the reframed narrative: dt-embedding is one channel through which the RSSM fidelity gap is exploited, not the only one.
As demonstrated by 008f's collapse without any dt-embedding, the biased RSSM state predictions themselves constitute an alternative exploitation channel.

%==============================================================================
\section{Discussion}\label{sec:discussion}
%==============================================================================

\subsection{The RSSM Fidelity Ceiling}

Our results converge on a single explanatory framework: the RSSM's prediction quality at $\dt>1$ is fundamentally insufficient to support value bootstrapping in imagination.
The fidelity gap (3.6$\times$ observation, 5.1$\times$ reward at $\dt=2$) creates a \emph{ceiling} that no conditioning mechanism can breach.

This ceiling manifests through two failure modes, depending on whether overestimation occurs:
\begin{enumerate}[nosep]
  \item \textbf{With overestimation} ($\hat{r} > 1.0$, subsample ON): RSSM learns biased $\dt=2$ dynamics $\to$ value bootstrapping amplifies bias $\to$ delusional critic $\to$ actor entropy increases $\to$ collapse.
    Speed depends on whether dt-embedding accelerates exploitation (008d: fast) or not (008f: delayed, same endpoint).
  \item \textbf{Without overestimation} ($\hat{r} < 1.0$, subsample OFF): RSSM never learns $\dt=2$ dynamics $\to$ $\dt=2$ imagination is out-of-distribution $\to$ value function treats it as noise $\to$ no collapse, but $\dt=2$ imagination provides no useful signal $\to$ slow degradation, reaching at best $\sim$60\% of baseline on average across seeds.
\end{enumerate}

Both paths lead to substantially degraded performance relative to $\dt=1$.
Temporal abstraction fails whether or not the world model has seen multi-step transitions.

\subsection{Why Single-Seed Results Are Misleading}

Our multi-seed analysis reveals that several phenomena initially presented as systematic mechanisms are in fact seed-specific artifacts:
\begin{itemize}[nosep]
  \item The V-shaped recovery in 008a (seed 0: 501$\to$778) does not replicate (seeds 1--2: 420, 480). Mean recovery is 59\% of baseline, not 85\%.
  \item 008d collapse occurs in 2/3 seeds, not 3/3. One seed survives in a degraded-but-stable state.
  \item 008f collapse (delayed) also occurs in 2/3 seeds---the \emph{same} rate as 008d---eliminating the dt-embedding causal argument entirely.
\end{itemize}

The cross-seed eval std at 45K for 008d is 75 (vs baseline std 18 asymptotically), a $>$4$\times$ increase in variance.
This variance inflation is itself a signature of operating near the RSSM fidelity ceiling: small differences in initialization determine whether the system falls into the collapse attractor or stabilizes in a degraded equilibrium.

\subsection{dt-Embedding: Accelerator, Not Cause}

Table~\ref{tab:hopper} summarizes the dt-embedding effect size across environments: 29\% on Walker Walk, 3\% on Cheetah Run, 5\% on Hopper Hop.
The large Walker Walk effect is environment-specific, likely reflecting Walker's particular sensitivity to dt-conditioning in its reward landscape.
In the other two environments, the effect is within seed variance.

The dt-embedding functions as an \emph{acceleration factor}: when present and trained, it provides a precise channel through which the value function can exploit the RSSM's biased predictions, hastening collapse by $\sim$15K steps.
When absent, the exploitation proceeds through the biased state predictions themselves---slower but equally inevitable.

\subsection{Implications for Temporal Abstraction in World-Model RL}

The RSSM fidelity ceiling has architectural roots: the GRU-based recurrence loses intermediate state information that cannot be recovered from the latent state alone.
Locomotion tasks exacerbate this because reward is a nonlinear function of contact events that occur at the skipped timesteps.

Future approaches to temporal abstraction in world-model RL must address the prediction quality bottleneck directly, rather than layering conditioning mechanisms atop an insufficiently accurate model.
Potential directions include:
\begin{itemize}[nosep]
  \item Architectures with explicit multi-scale temporal modeling (e.g., hierarchical state spaces)
  \item Uncertainty-aware imagination that downweights or truncates trajectories when prediction confidence drops
  \item Training objectives that explicitly penalize $\dt>1$ prediction error rather than treating all timestep predictions equally
\end{itemize}

\subsection{Methodological Contributions}

The factorial diagnostic methodology employed here---systematically crossing factors (dt-embedding type $\times$ subsampling $\times$ horizon) and validating across environments and seeds---provides a template for diagnosing failure modes in model-based RL more broadly.
The normalized v/e overshoot ratio ($\hat{r}$) and actor entropy serve as portable, early-warning diagnostics that can be computed from standard training logs without additional instrumentation.
These tools detected 008f's delayed collapse $\sim$10K steps before eval return began declining, and consistently flagged pathological regimes ($\hat{r} > 1.5$) across all environments and configurations.

%==============================================================================
\section{Conclusion}\label{sec:conclusion}
%==============================================================================

We have presented the first systematic, multi-environment, multi-seed empirical investigation of temporal abstraction in RSSM-based world models.
Our findings are unambiguous: temporal abstraction via variable-rate imagination is comprehensively broken in current latent-dynamics architectures.
No combination of dt-embedding configuration, subsampling strategy, or imagination horizon achieves competitive performance with the $\dt=1$ baseline.
The best $\dt=2$ configuration across three DMC environments reaches only 22--54\% of baseline performance on average across seeds.

The root cause is not the dt-embedding---an initial hypothesis that our extended experiments refute---but rather the RSSM's fundamental inability to predict multi-step transitions with sufficient accuracy.
The 3.6--5.1$\times$ prediction quality gap at $\dt=2$ establishes an \emph{RSSM fidelity ceiling} that cannot be overcome by conditioning mechanisms alone.
This fidelity gap propagates through value bootstrapping to produce either catastrophic collapse (when overestimation occurs) or persistent degradation (when it does not).

We contribute a factorial diagnostic methodology, the normalized v/e overshoot ratio as an early-warning metric, and the RSSM fidelity ceiling hypothesis as a framework for understanding why temporal abstraction fails and what must change for it to succeed.
Until world models can predict multi-step transitions with fidelity comparable to single-step predictions, temporal abstraction in imagination-based RL will remain fundamentally limited.

%==============================================================================
% References
%==============================================================================
\bibliographystyle{plainnat}

\begin{thebibliography}{10}

\bibitem[Hafner et~al.(2020)]{hafner2020dream}
D.~Hafner, T.~Lillicrap, J.~Ba, and M.~Norouzi.
\newblock Dream to control: Learning behaviors by latent imagination.
\newblock In \emph{ICLR}, 2020.

\bibitem[Hafner et~al.(2021)]{hafner2021mastering}
D.~Hafner, T.~Lillicrap, M.~Norouzi, and J.~Ba.
\newblock Mastering {A}tari with discrete world models.
\newblock In \emph{ICLR}, 2021.

\bibitem[Hafner et~al.(2023)]{hafner2023dreamerv3}
D.~Hafner, J.~Pasukonis, J.~Ba, and T.~Lillicrap.
\newblock Mastering diverse domains through world models.
\newblock \emph{arXiv:2301.04104}, 2023.

\bibitem[Sutton et~al.(1999)]{sutton1999options}
R.~S.~Sutton, D.~Precup, and S.~Singh.
\newblock Between {MDP}s and semi-{MDP}s: A framework for temporal abstraction in reinforcement learning.
\newblock \emph{Artificial Intelligence}, 112(1-2):181--211, 1999.

\bibitem[Sharma et~al.(2017)]{sharma2017figar}
S.~Sharma, A.~S.~Lakshminarayanan, and B.~Ravindran.
\newblock Fine grained action repetition for deep reinforcement learning.
\newblock In \emph{ICLR}, 2017.

\bibitem[Janner et~al.(2019)]{janner2019mbpo}
M.~Janner, J.~Fu, M.~Zhang, and S.~Levine.
\newblock When to trust your model: Model-based policy optimization.
\newblock In \emph{NeurIPS}, 2019.

\bibitem[Delalleau et~al.(2019)]{delalleau2019delusional}
O.~Delalleau, M.~Peter, E.~Bhatt, and D.~Khashabi.
\newblock Discrete-continuous mixtures in differentiable model-based reinforcement learning.
\newblock \emph{arXiv:1909.12830}, 2019.

\bibitem[Micheli et~al.(2023)]{micheli2023iris}
V.~Micheli, E.~Alonso, and F.~Fleuret.
\newblock Transformers are sample-efficient world models.
\newblock In \emph{ICLR}, 2023.

\bibitem[Robine et~al.(2023)]{robine2023twm}
J.~Robine, M.~H{\"o}ftmann, T.~Uelwer, and S.~Harmeling.
\newblock Transformer-based world models are happy with 100k interactions.
\newblock In \emph{ICLR}, 2023.

\end{thebibliography}

\end{document}
